\chapter{The Slip Becomes A Number}

\section{Where the loop fails to close}

The last chapter made measurement a loop and its reading the loop's alteration. This chapter finds the sharpest form of that
alteration: the place where the loop \emph{fails to close} at all --- where the returned state does not merely rotate by a
small residue but slips, discontinuously, from one regime to another. The failure is not a breakdown of the instrument. It is
its clearest signal. The instrument reads best exactly where the loop stops holding, because that is where the residue stops
being infinitesimal and becomes a magnitude the bench can name.

The physical image the construction keeps is the onset of slipping --- the transition from a circuit that sticks to one that
gives. Below a threshold the loop holds: parallel transport around it returns the state to itself, the holonomy is the
identity, nothing slips. Static. At the threshold the loop can no longer hold the accumulated residue, and the state slips
--- the transport returns rotated by a finite angle, the holonomy jumps away from the identity. Kinetic. The changeover is
abrupt because the two regimes are qualitatively different: one has a closed loop and no reading, the other has an open one
and a magnitude. The failure to close is the boundary between them, and the boundary is where the number is born. What was a
smooth accumulation of residue becomes, at the slip, a definite finite thing.

The geometry of the slip is the anholonomy: the net rotation a frame suffers when carried around a circuit that does not
close on the identity. Carry an orthonormal frame around the loop and it returns rotated by
\[
  \Delta\theta \;=\; \oint_\gamma \omega \;=\; \iint_{\Sigma} R ,
\]
the integral of the connection around the boundary, equal by Stokes to the curvature over the enclosed surface. When the
enclosed residue is below the threshold the return rotation is a correction; when it crosses the threshold the frame slips by
a full, readable angle. This is the slip made exact: $\Delta\theta$ is the failure of the loop to close, and it is finite,
signed, and --- this is the point --- a genuine magnitude rather than a mere distinction. The loop that closes gives zero;
the loop that slips gives $\Delta\theta$, and $\Delta\theta$ is the first number the instrument reads that is not a count.

Why the failure, and not the success, is the signal deserves stating plainly, because it inverts the naive expectation. One
might think an instrument works best when its loops close cleanly and returns nothing when they slip. The opposite is true. A
loop that closes has measured nothing --- it found no residue, reported no alteration, returned the identity. A loop that
slips has found the residue and made it readable. The slip is where the curvature stops hiding inside an infinitesimal
correction and stands out as a finite angle. So the instrument does not avoid the slip; it seeks it. It drives the loop
toward the threshold on purpose, because the threshold is the one place the residue is large enough to read and definite
enough to bracket.

In the two verbs, the slip is where a funge breaks and a tange becomes possible. Below threshold the loop funges its start
and end into one class --- ``no change'' --- and there is nothing to select. At the slip that funge fails: start and end fall
into different classes, the loop returns altered, and now there is a definite characteristic --- the rotation angle
$\Delta\theta$ --- to tange out and read. The slip is the failure of the covariant funge that first makes a contravariant
reading available. So the loop's failure to close is not the instrument breaking; it is the instrument switching on. The next
section takes the finite angle the slip produced and does the honest thing with it: brackets it, rather than guessing it.

\section{Bracketed, not guessed}

The slip produced a finite magnitude, $\Delta\theta$, the loop's failure to close. This section is about how the instrument
is permitted to report that magnitude, and the permission is strict: it may bracket the slip between two walls it can
compute, but it may never guess it, never fit it, never hand back a single confident point. The bracket is not a weaker
answer than the point. It is the honest one, and this section builds the discipline that will govern every reading in the
book, up to and including the coupling.

The mathematics is bounding from both sides. The slip is not known as an exact real --- recall that reals are only ever
approached, never completed --- so the instrument brackets it: it computes a lower wall $\ell$ and an upper wall $u$,
certifies $\ell < \Delta\theta < u$, and reports the interval $[\ell, u]$. The certification is a genuine one --- the walls
are constructed so that the true value provably lies between them, not so that it plausibly does. This is the difference
between a bound and a fit. A bound is a pair of walls with a proof that the value is inside; a fit is a single number chosen
to sit agreeably near the data, with no proof it is the value at all. The instrument produces bounds. It refuses fits. The
whole reason the reals were built as approached-never-completed is so that the honest report of one is an interval and not a
point.

The temptation the discipline refuses is the fitted value, and it is worth naming exactly, because it is the crank's
signature. Given a slip and a target, one can always choose a number that lands near the target and present it as the
reading. That number is not measured; it is fitted --- selected after the fact to match. It carries the appearance of
precision and none of the substance, because nothing certifies it but its agreeableness. The instrument is built so that this
step is not available to it: it does not select a value to match a target, it drives two walls toward each other and reports
where they are. There is no slot in the construction where a target is consulted, because a value chosen against a target is
exactly the falsification the whole method exists to refuse. Bracket, never guess: the two walls are computed from the
geometry, and the target is never in the room.

There is a reason the bracket is trustworthy in a way a fit never is, and it is the reason the whole book insists on
intervals. A bracket is \emph{falsifiable in both directions}: if the true value lay outside $[\ell, u]$ the certification
would fail, and the failure would be detectable. A fitted point makes no such promise --- it cannot be outside itself, so it
can never be caught being wrong. The interval carries its own honesty inside it: the width $u - \ell$ is an explicit
statement of what the instrument does not know, published alongside what it does. A point hides its ignorance; a bracket
declares it. So the bracket is not the instrument admitting weakness; it is the instrument publishing its resolution, which
is the one thing a fit can never do.

In the two verbs, bracketing is a funge the instrument closes only as far as its resolution allows, and refuses to tange
past. It funges the slip into the interval $[\ell, u]$ --- bags it between two computed walls --- and declines the final
tange that would select a single point from inside the bag. The point would be a tange the resolution cannot support; the
interval is the funge the resolution earns. So the slip is bracketed, not guessed, because the honest gathering is the widest
one the walls certify, and the dishonest one would be a selection the instrument has no ground to make. The next section
takes this bracketed slip and marks the moment it stops being a symbol and becomes, for the first time, an instrument
reading.

\section{The reading}

The slip has been made a finite magnitude and then bracketed between certified walls. This closing section marks what has
quietly happened across the chapter: the word ``number'' has changed meaning. It began the book as a symbol --- a label, a
position, a count of distinctions. It ends this chapter as an \emph{instrument reading} --- a bracketed, dimensional
magnitude, produced by a loop's failure to close and reported as an interval the bench can repeat. This is the hinge of the
whole construction, the place where the mathematics of distinction turns into the physics of measurement.

Take the change precisely. Earlier a number was the size of a funge, an integer count of told-apart states --- a symbol in
the sense that it lived on the labels and carried no scale. Now a number is a slip: a rotation angle $\Delta\theta$, bracketed
in $[\ell, u]$, carrying units, produced by a physical circuit, and reproducible on the bench. The two are not the same kind
of thing. The count answered \emph{how many}; the reading answers \emph{how much}, to a stated resolution, with a certified
interval. The reading has a dimension the count never had --- an angle, a length, an energy --- because it came from the
metric-laden loop and not from the bare distinction. Number, at this tier, has become a measured quantity, and a measured
quantity is a bracket with a unit.

The reading is a reading precisely because it is repeatable, and repeatability is the physics of the whole matter. Run the
loop again and it slips again by the same $\Delta\theta$, brackets again to the same $[\ell, u]$, because the residue it reads
is the invariant curvature that no relabeling alters. A quantity that returned a different interval each time you asked would
be no reading at all --- it would be noise wearing a number's clothes. The instrument reading is the interval that comes back
the same, and its coming-back-the-same is inherited straight from the invariance of the residue: the world's covariant
curvature is stable, so the bracket that reads it is stable. Repeatability is not an extra virtue bolted onto the reading; it
is the residue's invariance, seen through the loop.

This is also where the book's central refusal takes its mature form for the first time. The reading is an interval, and the
instrument will not narrow that interval to a point past the resolution that certifies it. It has a magnitude, and it holds
it as a bracket, because the bracket is what it can repeat and the point is what it would have to fit. Every later reading in
the book --- the faces, the field, the coupling itself --- is this same object: a dimensional magnitude, produced by a loop,
bracketed at a floor, and repeatable. The coupling at the climax will differ from this slip only in which loop produces it
and how deep the floor lies. The kind of thing it is --- a bracketed instrument reading, never a fitted point --- is fixed
here, at the first slip, and does not change.

In the two verbs, the reading is a funge the instrument owns and refuses to over-tange. It funges the slip into a repeatable
interval --- bags the magnitude between walls it can certify and rerun --- and declines to tange the interval down to the
single value a fit would announce. The reading is the funge it can repeat; the point would be the tange it cannot support. So
the slip becomes a number in exactly the sense the whole book means by ``number'' from here on: a bracketed, dimensional,
repeatable reading of an invariant residue, held open at the floor of resolution. What that floor \emph{is} --- how deep it
lies, why nothing registers below it, and why the reading is the interval and not the point --- is the object of the next
chapter, which builds the bracket into the instrument's permanent output.
